\documentclass[11pt]{article}
\usepackage[a4paper, top=18mm, bottom=18mm, left=20mm, right=20mm]{geometry}
\usepackage[T2A]{fontenc}
\usepackage[utf8]{inputenc}
\usepackage[ukrainian]{babel}
\usepackage{graphicx}
\usepackage[export]{adjustbox}
\usepackage{amsmath}
\usepackage{amssymb}
\usepackage{microtype}
\graphicspath{{/app/Quiz/Scripts/2023/ProgAll/15BinTree/}}
\setlength{\parindent}{0pt}
\setlength{\parskip}{6pt}
\emergencystretch=3em
\pagestyle{empty}
\begin{document}
%begin
{\centering\Large\bfseries Контрольна робота\par}
\medskip
\noindent\rule{\linewidth}{0.5pt}
\smallskip
\noindent\textbf{Варіант \No\,1}\hfill \textit{ПІБ:}\,\underline{\hspace{60mm}}
\vspace{4mm}

%beginitem
1. %goodpath:Scripts/2018/Mod2018_1/03-good.txt
\textbf{Відмітити все, що є коректним оператором С++:}
( 1 ) double x=1;

( 2 ) cin<<3;

( 3 ) x=y=z

( 4 ) cin>>x+3;
\vspace{5mm}
%enditem

%beginitem
2. \textbf{Що буде виведено за виконання фрагменту коду?}

%code
cout << (5 % 5 % 5 % 5);
%code
\vspace{5mm}
%enditem

%beginitem
3. \textbf{Що буде виведено за виконання фрагменту коду?}

a,b,c=1,9,4
def h(b):
global c    a*=3
    b=1
    c=5
    return a+b+c

a,b,c=5,7,3
print(h(a),a,b,c)
\vspace{5mm}
%enditem

%beginitem
4. \textbf{Що буде виведено за виконання фрагменту коду?}

%code
print(3.0 > 3 is 4 == False)
%code
\vspace{5mm}
%enditem

%beginitem
5. \textbf{Що буде виведено за виконання фрагменту коду?}
%code
a, b, c = 7, 8, 6
def h(a, b, c=9):
    print(a, b, c, end=" ")

h(1, 2, 4)
print(a, b, c)
%code
\vspace{5mm}
%enditem

%beginitem
6. %goodpath:Scripts/2019/Mod2019_1/Lex/01-good.txt
\textbf{Відмітити все, що є однією правильною лексемою:}
( 1 ) true

( 2 ) 1+2j

( 3 ) xy

( 4 ) 1+1j
\vspace{5mm}
%enditem

%beginitem
7. \textbf{Що буде виведено за виконання фрагменту коду?}

print('R', end=' ')
b=['0','1','2','3','4']
print(b.pop(-4), end=' ')
print(b)

\vspace{5mm}
%enditem

%beginitem
8. \textbf{Що буде виведено за виконання фрагменту коду?}

%code
4**-3--1
%code
\vspace{5mm}
%enditem

%beginitem
9. 

\begin{center}
	\adjustimage{max size={0.8\linewidth}{0.8\paperheight}}
	{../15BinTree/trees_exp/52.png}
\end{center}
Указати послідовність міток вершин дерева, зображеного на рис., що утвориться в результаті обходу дерева в оберненому порядку. Мітки записувати через пробіл.\par Алгоритм починає роботу з кореня (на рис. це верхня вершина).
%answer:52 оберненому
\vspace{5mm}
%enditem

%beginitem
10. %goodpath:Scripts/2019/Mod2019_1/Lex/03-good.txt
\textbf{Відмітити все, що є коректним оператором С++:}
( 1 ) 3**-1

( 2 ) if ('x'><'y'): a=6

( 3 ) x<y<z

( 4 ) (2+2j)%2
\vspace{5mm}
%enditem

%end
\newpage
\end{document}

